\subsection{Solving Equations with Two Exponentials} 
We have already seen a method for solving an equation where two exponential expressions have the same base. If the bases are different, we can use the logarithm in any base. (We usually choose the \makeblank[1in] log.)
\begin{example}
Solve $3^{x+5}=7^{2x-1}$
\begin{lpic}{}{11}
\end{lpic}
\end{example}
Note that exponential expressions are always \makeblank[1in], so we can easily have an equation with no solutions:
\begin{example}
Solve $3^{x+4}+11^{x-3}=0$
\begin{itemize}
\item $\makeblank[1in]>0$ for every value of $x$
\item $\makeblank[1in]>0$ for every value of $x$, too.
\item So $\makeblank[1in]+\makeblank[1in]>0$ for every value of $x$.
\end{itemize}
Since the expression on the the left is always \makeblank[1in] than 0, it can never \makeblank[1in] 0, so the equation has no solution.
\end{example}
We also can't deal with addition (in general).
\begin{example}
Solve $4^x=3^x+2$
\begin{lpic}{}{2}
\end{lpic}
We're stuck, because we can't split up a log of a \makeblank[1in]. For an equation like this, we will usually need to resort to \makeblank[1in] techniques.
\end{example}
Usually, an exponential equation with more than \makeblanksmall terms cannot be solved.